\documentclass[11pt,a4paper]{article}
% =====================================================================
%  Gemeinsame Vorlage der Arbeitsblätter – 1. Formale Sprachen (Q13)
%  (gleiches Layout wie informatik/13/material/grundwissen/ab-vorlage.tex
%  und informatik/12/rekursive-datenstrukturen/arbeitsblaetter/ab-vorlage.tex)
%  Einbinden in jedem Blatt direkt nach \documentclass: % =====================================================================
%  Gemeinsame Vorlage der Arbeitsblätter – 1. Formale Sprachen (Q13)
%  (gleiches Layout wie informatik/13/material/grundwissen/ab-vorlage.tex
%  und informatik/12/rekursive-datenstrukturen/arbeitsblaetter/ab-vorlage.tex)
%  Einbinden in jedem Blatt direkt nach \documentclass: % =====================================================================
%  Gemeinsame Vorlage der Arbeitsblätter – 1. Formale Sprachen (Q13)
%  (gleiches Layout wie informatik/13/material/grundwissen/ab-vorlage.tex
%  und informatik/12/rekursive-datenstrukturen/arbeitsblaetter/ab-vorlage.tex)
%  Einbinden in jedem Blatt direkt nach \documentclass: \input{ab-vorlage}
%  Übersetzen mit XeLaTeX, LuaLaTeX, pdfLaTeX oder Tectonic.
%
%  Blatt:   tectonic ab-1-2-formale-sprachen.tex
%  Lösung:  tectonic ab-1-2-formale-sprachen-lsg.tex
%           (Datei enthält nur: \def\mitloesung{}\input{ab-1-2-formale-sprachen};
%            füllt alle Lücken rot aus; siehe \lsg, \lsglinien, \lsgoder, lsgbild, \karolsg)
% =====================================================================
\usepackage[a4paper,left=16mm,right=16mm,top=23mm,bottom=13mm,headheight=12mm,headsep=3mm,footskip=6mm]{geometry}
\usepackage{iftex}
\ifPDFTeX
  \usepackage[T1]{fontenc}
  \usepackage[utf8]{inputenc}
  \usepackage{helvet}
  \usepackage[scaled=0.86]{beramono}
  \renewcommand{\familydefault}{\sfdefault}
\else
  \usepackage{fontspec}
  \setmainfont{texgyreheros}[Extension=.otf,UprightFont=*-regular,BoldFont=*-bold,ItalicFont=*-italic,BoldItalicFont=*-bolditalic]
  \setmonofont{DejaVuSansMono}[Extension=.ttf,UprightFont=*,BoldFont=*-Bold,ItalicFont=*-Oblique,BoldItalicFont=*-BoldOblique,Scale=0.86]
\fi
\usepackage[ngerman,shorthands=off]{babel}
\usepackage{amssymb}
\usepackage{xcolor,graphicx,tabularx,array,enumitem,fancyhdr,tikz,multicol,calc,etoolbox,pifont}
\usepackage[most]{tcolorbox}
\usetikzlibrary{arrows.meta,positioning,calc}
\usepackage[hidelinks]{hyperref}
\setlength{\parindent}{0pt}
\setlength{\parskip}{3pt}
\setlist{nosep,leftmargin=*}

% ---------- Lösungsschalter ----------
\newif\ifloesung
\ifdefined\mitloesung\loesungtrue\fi
\newcommand{\lsgfarbe}{\color{red!75!black}}

% ---------- Kopf- und Fußzeile (einheitliche Form) ----------
\newcommand{\lehrkraft}{Hau}
\newcommand{\themenbereich}{1. Formale Sprachen}
\newcommand{\abnummer}{}
\pagestyle{fancy}\fancyhf{}
\renewcommand{\headrulewidth}{0.4pt}
\fancyhead[L]{\small Klasse 13\\Datum:}
\fancyhead[C]{\small Themenbereich:\\\themenbereich}
\fancyhead[R]{\small \lehrkraft\ifloesung\\{\lsgfarbe\bfseries Lösung}\fi}
\fancyfoot[C]{\scriptsize Informatik 13 gA\,·\,Blatt \abnummer\ifloesung\ (Lösung)\fi\,·\,Seite \thepage}
\newcommand{\abtitel}[2]{\renewcommand{\abnummer}{#1}\begin{center}\Large\bfseries #1\quad\underline{#2}\end{center}\vspace{-2mm}}

% ---------- Boxen ----------
\newenvironment{aufgabe}[2][]{\begin{tcolorbox}[enhanced,breakable,colback=white,colframe=black,boxrule=0.7pt,arc=3mm,left=2mm,right=2mm,top=1mm,bottom=1.2mm,before skip=2mm,after skip=2mm]\textbf{Aufgabe #2)}\ifstrempty{#1}{}{ \textit{#1}}\par\vspace{0.3mm}}{\end{tcolorbox}}
\newtcolorbox{merke}[1][Merke]{enhanced,breakable,colback=green!5,colframe=green!40!black,boxrule=0.8pt,arc=2mm,left=2mm,right=2mm,top=1mm,bottom=1.2mm,before skip=2mm,after skip=2mm,title={#1},fonttitle=\bfseries,coltitle=white,toptitle=0.3mm,bottomtitle=0.3mm}
\newtcolorbox{hinweis}{enhanced,breakable,colback=yellow!12,colframe=orange!70!black,boxrule=0.6pt,arc=2mm,left=2mm,right=2mm,top=1mm,bottom=1mm,before skip=2mm,after skip=2mm}

% ---------- Lücken, Linien, Karos ----------
\newcommand{\luecke}[1][3.5cm]{\rule[-0.5ex]{#1}{0.4pt}}
\newcommand{\linien}[1]{\par\vspace{1.5mm}\foreach \i in {1,...,#1}{\noindent\rule[0pt]{\linewidth}{0.3pt}\par\vspace{5.5mm}}\vspace{-3mm}}
\newcommand{\code}[1]{\texttt{#1}}
\newcommand{\eps}{$\varepsilon$}

% Lücke mit Lösung: \lsg[Mindestbreite]{Lösungstext}
%  Die Lösung steht (ohne eigene Höhe/Tiefe) auf derselben Linie wie im Blatt: Das Layout verschiebt sich nicht,
%  solange der Text nicht breiter als die Mindestbreite ist.
\newlength{\lsgbreite}
\newcommand{\lsg}[2][3.5cm]{\leavevmode\ifloesung\settowidth{\lsgbreite}{\,#2\,}\ifdim\lsgbreite<#1\setlength{\lsgbreite}{#1}\fi\rlap{\smash{\raisebox{0.8pt}{\makebox[\lsgbreite][l]{\lsgfarbe\,#2}}}}\luecke[\lsgbreite]\else\luecke[#1]\fi}
% Schreiblinien mit Lösung: \lsglinien{Anzahl}{Lösungstext} (gleiche Höhe wie \linien:
%  je Linie \baselineskip + 5,5 mm + \parskip, ausgemessen)
\newlength{\lsglinienhoehe}
\newcommand{\lsglinien}[2]{\ifloesung\par\vspace{1.5mm}\setlength{\lsglinienhoehe}{\dimexpr #1\dimexpr\baselineskip+5.5mm+\parskip\relax-16.25pt\relax}\noindent\begin{minipage}[t][\lsglinienhoehe][t]{\linewidth}\lsgfarbe #2\end{minipage}\par\vspace{-1.5mm}\else\linien{#1}\fi}
% Beliebiger Inhalt: \lsgoder{nur in der Lösung}{nur im Blatt}
\newcommand{\lsgoder}[2]{\ifloesung\leavevmode#1\else#2\fi}
% Text nur in der Lösung sichtbar, belegt im Blatt aber denselben Platz
\newcommand{\lsgtext}[1]{\ifloesung{\lsgfarbe#1}\else\phantom{#1}\fi}
% Ankreuzfeld: \kreuz{1} = angekreuzt (nur in der Lösung), \kreuz{0} = leer
\newcommand{\kreuz}[1]{\leavevmode\ifloesung\ifnum#1=1 {\lsgfarbe$\boxtimes$}\else$\square$\fi\else$\square$\fi}

% ---------- Syntaxdiagramme (Maße in mm) ----------
%  Terminal rund (orange), Nichtterminal eckig (grün), Linien mit Pfeilen:
%  \draw[sdl,->] (0,0) -- ++(6,0) node[nt,anchor=west] (a) {Augen};   (Pfeil endet am Kasten)
\tikzset{
  sd/.style={x=1mm,y=1mm,>={Stealth[length=1.6mm,width=1.3mm]},font=\small},
  sdl/.style={line width=0.55pt,rounded corners=1.2mm},
  te/.style={draw,rounded corners=2.1mm,minimum height=4.2mm,minimum width=4.2mm,inner sep=1pt,font=\ttfamily\small,fill=orange!18},
  nt/.style={draw,minimum height=4.2mm,inner sep=1.6pt,font=\small,fill=green!12},
  sdname/.style={anchor=south west,font=\small\bfseries,inner sep=0pt},
  punkte/.style={font=\small,inner sep=0pt},
  % Ableitungsbaum
  baum/.style={font=\small,edge from parent/.style={draw,line width=0.5pt,-{Stealth[length=1.4mm]}},level distance=8.5mm},
}
% Lösungszeichnung: im Blatt unsichtbar, belegt aber denselben Platz
\ifloesung
  \tikzset{lsgbild/.style={red!75!black,te/.append style={fill=red!5},nt/.append style={fill=red!5}}}
\else
  \tikzset{lsgbild/.style={draw opacity=0,fill opacity=0,text opacity=0}}
\fi
% Karofläche, in der Lösung mit roter Zeichnung: \karolsg{Anzahl Karozeilen}{TikZ-Code}
% Koordinaten der Zeichnung in mm, Ursprung links oben, y nach unten negativ.
\newcommand{\karolsg}[2]{\par\vspace{1mm}\noindent\begin{tikzpicture}
  \draw[step=5mm,gray!45,thin] (0,0) grid (\linewidth-1pt,#1*5mm);
  \useasboundingbox (0,0) rectangle (\linewidth-1pt,#1*5mm);
  \ifloesung\begin{scope}[sd,shift={(0,#1*5mm)},lsgbild]#2\end{scope}\fi
\end{tikzpicture}\par}

% ---------- Bausteine für Syntaxdiagramme (innerhalb tikzpicture[sd]) ----------
%  Jeder Baustein beginnt an einem Punkt (Knoten-Anker oder Koordinate) und
%  legt die Koordinate <name>-o an, an der es weitergeht.
% Aneinanderreihung: \sdnext{name}{von}{te|nt}{Text}   (weiter bei name.east)
\newcommand{\sdnext}[4]{\draw[sdl,->] (#2) -- ++(4,0) node[#3,anchor=west] (#1) {#4};\coordinate (#1-o) at (#1.east);}
% Option: \sdopt{name}{von}{te|nt}{Text}
\newcommand{\sdopt}[4]{\path (#2) ++(4,-5.5) node[#3,anchor=west] (#1) {#4};
  \coordinate (#1-x) at ($(#1.east)+(3.5,0)$); \coordinate (#1-o) at (#1-x |- #2);
  \draw[sdl,->] (#2) |- (#1.west); \draw[sdl,->] (#1.east) -| (#1-o); \draw[sdl] (#2) -- (#1-o);}
% Wiederholung, mindestens einmal: \sdwdh{name}{von}{te|nt}{Text}
\newcommand{\sdwdh}[4]{\draw[sdl,->] (#2) -- ++(6,0) node[#3,anchor=west] (#1) {#4};
  \coordinate (#1-o) at ($(#1.east)+(4,0)$);
  \draw[sdl] (#1.east) -- (#1-o); \draw[sdl,->] (#1-o) -- ++(0,-5.5) -| ($(#2)+(2.5,0)$);}
% Wiederholung, beliebig oft (auch keinmal): \sdwdhn{name}{von}{te|nt}{Text}
\newcommand{\sdwdhn}[4]{\path (#2) ++(6,-5.5) node[#3,anchor=west] (#1) {#4};
  \coordinate (#1-x) at ($(#1.east)+(4,0)$); \coordinate (#1-o) at (#1-x |- #2);
  \draw[sdl] (#2) -- (#1-o); \draw[sdl,->] (#1-o) |- (#1.east); \draw[sdl,->] (#1.west) -| ($(#2)+(2.5,0)$);}
% Alternative: \sdalt{name}{von}{Breite bis zum Zusammenführen}{Zeilenabstand}{te/a, nt/Ziffer, pk/, ...}
%  pk/ setzt senkrechte Pünktchen (…) statt eines Zweigs
\newcommand{\sdalt}[5]{\coordinate (#1-o) at ($(#2)+(#3,0)$);
  \foreach \typ/\txt [count=\i from 0] in {#5}{%
    \ifdefstring{\typ}{pk}{\node[punkte] at ($(#2)+(6.1,-\i*#4+0.4)$) {$\vdots$};}{%
      \ifdefstring{\typ}{te}{\tikzset{sdakt/.style={te}}}{\tikzset{sdakt/.style={nt}}}%
      \node[sdakt,anchor=west] (#1\i) at ($(#2)+(4,-\i*#4)$) {\txt};
      \ifnum\i=0\relax \draw[sdl,->] (#2) -- (#1\i.west); \draw[sdl] (#1\i.east) -- (#1-o);
      \else \draw[sdl,->] (#2) |- (#1\i.west); \draw[sdl,->] (#1\i.east) -| (#1-o);\fi}}}
% Ende eines Diagramms
\newcommand{\sdende}[1]{\draw[sdl,->] (#1) -- ++(5,0);}

% ---------- Die Smiley-Grammatik (Blatt 1.3 und 1.4) ----------
\newcommand{\sdSmileyfolge}{\begin{tikzpicture}[sd,baseline=0mm]
  \node[sdname] at (0,3.3) {Smileyfolge}; \coordinate (s) at (0,0);
  \sdwdh{sm}{s}{nt}{Smiley}\sdende{sm-o}\end{tikzpicture}}
\newcommand{\sdSmiley}{\begin{tikzpicture}[sd,baseline=0mm]
  \node[sdname] at (0,3.3) {Smiley}; \coordinate (s) at (0,0);
  \sdnext{au}{s}{nt}{Augen}\sdnext{mi}{au-o}{nt}{Mitte}\sdnext{mu}{mi-o}{nt}{Mund}\sdende{mu-o}\end{tikzpicture}}
\newcommand{\sdAugen}{\begin{tikzpicture}[sd,baseline=0mm]
  \node[sdname] at (0,3.3) {Augen}; \coordinate (s) at (0,0);
  \sdalt{a}{s}{12}{5.5}{te/:, te/;}\sdende{a-o}\end{tikzpicture}}
\newcommand{\sdMitte}{\begin{tikzpicture}[sd,baseline=0mm]
  \node[sdname] at (0,3.3) {Mitte}; \coordinate (s) at (0,0);
  \sdopt{n}{s}{te}{-}\sdende{n-o}\end{tikzpicture}}
\newcommand{\sdMund}{\begin{tikzpicture}[sd,baseline=0mm]
  \node[sdname] at (0,3.3) {Mund}; \coordinate (s) at (0,0);
  \sdalt{m}{s}{12}{5.5}{te/), te/{(}, te/D, te/P, te/*}\sdende{m-o}\end{tikzpicture}}

%  Übersetzen mit XeLaTeX, LuaLaTeX, pdfLaTeX oder Tectonic.
%
%  Blatt:   tectonic ab-1-2-formale-sprachen.tex
%  Lösung:  tectonic ab-1-2-formale-sprachen-lsg.tex
%           (Datei enthält nur: \def\mitloesung{}\documentclass[11pt,a4paper]{article}
\input{ab-vorlage}
% Blatt:   tectonic ab-1-2-formale-sprachen.tex
% Lösung:  tectonic ab-1-2-formale-sprachen-lsg.tex
% Progression: Alphabet -> Wort/Länge -> Wortmenge A* (zählen: unendlich) -> Sprache L ⊆ A*
%              -> Transfer (Alltagssprachen, Syntax vs. Semantik) -> Smileys in Worten (mehrdeutig)
\newcommand{\zeile}{\rule[-2.2mm]{0pt}{7mm}}
\newcommand{\hoch}{\rule[-1.6mm]{0pt}{6mm}}
\begin{document}
\abtitel{1.2}{Formale Sprachen}

Ob \code{2 + * 5} oder \code{34.07.2026} korrekt gebildet sind, kann ein Computer nur prüfen, wenn \textbf{eindeutig} festgelegt ist, welche \emph{Zeichen} erlaubt sind und welche \emph{Zeichenfolgen} dazugehören.

\begin{merke}[Alphabet, Wort und Länge]
\begin{itemize}
\item Ein \textbf{Alphabet} $A$ ist eine \lsg[2.2cm]{endliche}, nichtleere Menge von \lsg[2cm]{Zeichen}.\\[0.5mm]
$A_1 = \{\code{0}, \code{1}\}$\qquad $A_2 = \{\code{a}, \code{b}, \code{c}, \dots, \code{z}\}$\qquad $A_3 = \{\code{:}, \code{;}, \code{-}, \code{)}, \code{(}, \code{D}, \code{P}, \code{*}\}$
\item Ein \textbf{Wort} (eine \textbf{Zeichenkette}) über $A$ ist eine endliche Folge von Zeichen aus $A$.\\
Das \textbf{leere Wort} \eps{} enthält \lsg[2.6cm]{kein Zeichen}.
\item Die \textbf{Länge} $|w|$ eines Wortes $w$ ist die \lsg[4cm]{Anzahl seiner Zeichen}. Beispiel: $|\code{0110}| = 4$, \ $|\varepsilon| = $ \lsg[0.8cm]{0}.
\end{itemize}
\end{merke}

\begin{aufgabe}{1}
\textbf{a)} Gib je drei Wörter unterschiedlicher Länge an und notiere ihre Länge.\\[1mm]
\renewcommand{\arraystretch}{1.1}
\begin{tabularx}{\linewidth}{|p{1.3cm}|X|X|X|}\hline
über $A_1$\zeile & $k_1 = \varepsilon$, \quad $|k_1| = 0$ & \lsgoder{\lsgfarbe $k_2 = \code{1}$, \quad $|k_2| = 1$}{} & \lsgoder{\lsgfarbe $k_3 = \code{0110}$, \quad $|k_3| = 4$}{}\\\hline
über $A_2$\zeile & \lsgoder{\lsgfarbe \code{x}, \quad Länge 1}{} & \lsgoder{\lsgfarbe \code{baum}, \quad Länge 4}{} & \lsgoder{\lsgfarbe \code{informatik}, \quad Länge 10}{}\\\hline
über $A_3$\zeile & \lsgoder{\lsgfarbe \code{;D}, \quad Länge 2}{} & \lsgoder{\lsgfarbe \code{:-)}, \quad Länge 3}{} & \lsgoder{\lsgfarbe \code{DD:*(}, \quad Länge 5}{}\\\hline
\end{tabularx}\\[1.2mm]
\textbf{b)} Ist \code{ab} ein Wort über $A_1$? Ist \code{:-)} ein Wort über $A_2$? \lsg[7.4cm]{Nein, beides nicht: Zeichen fehlen im Alphabet.}\\[1mm]
\textbf{c)} Es gibt genau ein Wort, das über \emph{jedem} Alphabet gebildet werden kann. Welches? Warum? \\[0.6mm]
\lsg[\linewidth]{Das leere Wort \eps: Es enthält kein Zeichen, also auch keines, das im Alphabet fehlen könnte.}\\[1mm]
\textbf{d)} Kreuze an, was richtig ist.\qquad
\kreuz{0} $\varepsilon \in A_1$\qquad \kreuz{1} $\code{0} \in A_1$\qquad \kreuz{1} \code{000} ist Wort über $A_1$\qquad \kreuz{0} $|\code{:-)}| = 1$
\end{aufgabe}

\begin{aufgabe}[Wie viele Wörter gibt es?]{2}
\textbf{a)} Schreibe \emph{alle} Wörter über $A_1$ der Länge 0 bis 3 auf und zähle sie. Zähle dann auch für $A_2$ und $A_3$.\\[1mm]
\renewcommand{\arraystretch}{1.1}
\begin{tabularx}{\linewidth}{|p{2.9cm}|p{0.85cm}|p{1.5cm}|p{2.8cm}|X|p{1.4cm}|}\hline
\textbf{Länge $n$} & \textbf{0} & \textbf{1} & \textbf{2} & \textbf{3} & \textbf{$n$}\\\hline
Wörter über $A_1$\rule[-6mm]{0pt}{11mm} & \lsgoder{\lsgfarbe \eps}{} & \lsgoder{\lsgfarbe \code{0}, \code{1}}{} & \lsgoder{\lsgfarbe \code{00}, \code{01}, \code{10}, \code{11}}{} & \lsgoder{\lsgfarbe \code{000}, \code{001}, \code{010}, \code{011}, \code{100}, \code{101}, \code{110}, \code{111}}{} & \\\hline
Anzahl über $A_1$\hoch & \lsgoder{\lsgfarbe 1}{} & \lsgoder{\lsgfarbe 2}{} & \lsgoder{\lsgfarbe 4}{} & \lsgoder{\lsgfarbe 8}{} & \lsgoder{\lsgfarbe $2^n$}{}\\\hline
Anzahl über $A_2$\hoch & \lsgoder{\lsgfarbe 1}{} & \lsgoder{\lsgfarbe 26}{} & \lsgoder{\lsgfarbe $26^2 = 676$}{} & \lsgoder{\lsgfarbe $26^3 = 17\,576$}{} & \lsgoder{\lsgfarbe $26^n$}{}\\\hline
Anzahl über $A_3$\hoch & \lsgoder{\lsgfarbe 1}{} & \lsgoder{\lsgfarbe 8}{} & \lsgoder{\lsgfarbe 64}{} & \lsgoder{\lsgfarbe 512}{} & \lsgoder{\lsgfarbe $8^n$}{}\\\hline
\end{tabularx}\\[1.2mm]
\textbf{b)} Wie viele Wörter gibt es \emph{insgesamt} über $A_1$? Begründe, obwohl $A_1$ nur zwei Zeichen hat.\\[0.6mm]
\lsg[\linewidth]{Unendlich viele: Zu jeder Länge $n$ gibt es $2^n$ Wörter, und die Länge ist nicht begrenzt.}
\end{aufgabe}

\begin{merke}[Wortmenge und formale Sprache]
\begin{itemize}
\item Die \textbf{Wortmenge} $A^*$ ist die Menge \lsg[2.6cm]{aller Wörter} über $A$, einschließlich \eps. Obwohl $A$ endlich ist, ist $A^*$ immer \lsg[2.4cm]{unendlich}.\\[0.5mm]
Beispiel: $A_1^* = \{$\,\lsg[9cm]{\eps, \code{0}, \code{1}, \code{00}, \code{01}, \code{10}, \code{11}, \code{000}, \code{001}, \dots}\,$\}$
\item Eine \textbf{formale Sprache} $L$ über $A$ ist eine \textbf{eindeutig festgelegte} Menge von Wörtern über $A$, also eine \lsg[2.4cm]{Teilmenge} von $A^*$: $L \subseteq A^*$. Eindeutig heißt: Für \emph{jedes} Wort aus $A^*$ steht fest, ob es zu $L$ gehört oder nicht.
\end{itemize}
\end{merke}

\newpage
\begin{aufgabe}[Sprachen über $A_2$]{3}
\begin{minipage}[t]{0.585\linewidth}
\vspace{0pt}
$L_1$ = alle Wörter aus $A_2^*$ mit Länge kleiner gleich 4\\
$L_2$ = alle Wörter aus $A_2^*$, die mit \code{b} beginnen\\
$L_3$ = alle Wörter aus $A_2^*$, die das Teilwort \code{inf} enthalten\\[1mm]
\textbf{a)} Kreuze an, zu welcher Sprache das Wort gehört.\\[1mm]
\textbf{b)} Welche der drei Sprachen sind endlich? \lsg[1.8cm]{nur $L_1$}\\[1mm]
\textbf{c)} Finde ein Wort, das zu $L_1$, $L_2$ \emph{und} $L_3$ gehört:\\[0.6mm] \lsg[3.2cm]{\code{binf}}\ \ Gibt es mehrere? \lsg[2.6cm]{nein, nur \code{binf}}
\end{minipage}\hfill
\begin{minipage}[t]{0.39\linewidth}
\vspace{0pt}
\renewcommand{\arraystretch}{1.12}
\begin{tabularx}{\linewidth}{|X|c|c|c|}\hline
\textbf{Wort} & $L_1$ & $L_2$ & $L_3$\\\hline
\code{baum} & \kreuz{1} & \kreuz{1} & \kreuz{0}\\\hline
\code{informatik} & \kreuz{0} & \kreuz{0} & \kreuz{1}\\\hline
\eps & \kreuz{1} & \kreuz{0} & \kreuz{0}\\\hline
\code{inf} & \kreuz{1} & \kreuz{0} & \kreuz{1}\\\hline
\code{binfo} & \kreuz{0} & \kreuz{1} & \kreuz{1}\\\hline
\end{tabularx}
\end{minipage}
\end{aufgabe}

\begin{aufgabe}[Formale Sprachen im Alltag]{4}
\textbf{a)} Ergänze die Tabelle.\\[1mm]
\renewcommand{\arraystretch}{1.1}
\begin{tabularx}{\linewidth}{|p{5.3cm}|X|p{2.4cm}|p{4.5cm}|}\hline
\textbf{Sprache} & \textbf{Alphabet} & \textbf{Wort aus $L$} & \textbf{Wort aus $A^*$, nicht aus $L$}\\\hline
Postleitzahlen\zeile & \lsgoder{\lsgfarbe \{0, 1, \dots, 9\}}{} & \code{83278} & \lsgoder{\lsgfarbe \code{824377} (6 Ziffern)}{}\\\hline
Binärzahlen ohne führende Null\zeile & \lsgoder{\lsgfarbe \{0, 1\}}{} & \lsgoder{\lsgfarbe \code{1011}}{} & \lsgoder{\lsgfarbe \code{0011}}{}\\\hline
Datum (TT.MM.JJJJ)\zeile & \lsgoder{\lsgfarbe \{0, 1, \dots, 9, .\}}{} & \lsgoder{\lsgfarbe \code{29.09.2026}}{} & \lsgoder{\lsgfarbe \code{29..092026}}{}\\\hline
\end{tabularx}\\[1.2mm]
\textbf{b)} Die Postleitzahl \code{99999} ist nicht vergeben. Gehört sie trotzdem zur Sprache „alle Wörter aus fünf Ziffern“? Was folgt daraus für Syntax und Semantik?
\lsglinien{2}{Ja, \code{99999} ist syntaktisch korrekt. Eine formale Sprache legt nur die \textbf{Syntax} fest; ob das Wort auch etwas bedeutet (\textbf{Semantik}), prüft sie nicht.}
\end{aufgabe}

\begin{aufgabe}[Die Sprache der Smileys]{5}
Im Chat werden Smileys aus den Zeichen von $A_3$ gebildet. Sie bilden eine Sprache $L_\text{Smiley} \subseteq A_3^*$.\\[1mm]
\textbf{a)} Kreuze an, welche Wörter du als Smiley erkennst.\\[1mm]
\kreuz{1} \code{:-)}\qquad \kreuz{1} \code{;D}\qquad \kreuz{1} \code{:-P}\qquad \kreuz{0} \code{)-:}\qquad \kreuz{0} \code{:-}\qquad \kreuz{1} \code{;*}\qquad \kreuz{0} \code{:--)}\qquad \kreuz{0} \code{DP}\qquad \kreuz{1} \code{:(}\\[1mm]
\textbf{b)} Beschreibe in Worten möglichst genau, welche Wörter aus $A_3^*$ zu $L_\text{Smiley}$ gehören.
\lsglinien{3}{Ein Smiley beginnt mit Augen (\code{:} oder \code{;}). Danach kann eine Nase \code{-} folgen, höchstens eine. Am Ende steht genau ein Mund: \code{)}, \code{(}, \code{D}, \code{P} oder \code{*}. Sonst folgt nichts.}
\textbf{c)} Vergleiche mit deinem Nachbarn: Gibt es Wörter, bei denen eure Beschreibungen verschieden entscheiden oder gar nicht entscheiden? Notiere ein Beispiel und woran es liegt.
\lsglinien{2}{z.\,B. \code{)-:} (gespiegelt?), \code{:--)} (wie viele Nasen?), \code{:-):-)} (mehrere Smileys?). Umgangssprache ist ungenau und mehrdeutig; man übersieht leicht Fälle.}
\end{aufgabe}

\begin{merke}[Erkenntnis]
Eine Beschreibung in Umgangssprache legt eine Sprache oft \lsg[2.8cm]{nicht eindeutig} fest. Nötig sind genaue \textbf{Regeln}, nach denen man für jedes Wort entscheiden kann, ob es zur Sprache gehört.
\end{merke}
\end{document}
;
%            füllt alle Lücken rot aus; siehe \lsg, \lsglinien, \lsgoder, lsgbild, \karolsg)
% =====================================================================
\usepackage[a4paper,left=16mm,right=16mm,top=23mm,bottom=13mm,headheight=12mm,headsep=3mm,footskip=6mm]{geometry}
\usepackage{iftex}
\ifPDFTeX
  \usepackage[T1]{fontenc}
  \usepackage[utf8]{inputenc}
  \usepackage{helvet}
  \usepackage[scaled=0.86]{beramono}
  \renewcommand{\familydefault}{\sfdefault}
\else
  \usepackage{fontspec}
  \setmainfont{texgyreheros}[Extension=.otf,UprightFont=*-regular,BoldFont=*-bold,ItalicFont=*-italic,BoldItalicFont=*-bolditalic]
  \setmonofont{DejaVuSansMono}[Extension=.ttf,UprightFont=*,BoldFont=*-Bold,ItalicFont=*-Oblique,BoldItalicFont=*-BoldOblique,Scale=0.86]
\fi
\usepackage[ngerman,shorthands=off]{babel}
\usepackage{amssymb}
\usepackage{xcolor,graphicx,tabularx,array,enumitem,fancyhdr,tikz,multicol,calc,etoolbox,pifont}
\usepackage[most]{tcolorbox}
\usetikzlibrary{arrows.meta,positioning,calc}
\usepackage[hidelinks]{hyperref}
\setlength{\parindent}{0pt}
\setlength{\parskip}{3pt}
\setlist{nosep,leftmargin=*}

% ---------- Lösungsschalter ----------
\newif\ifloesung
\ifdefined\mitloesung\loesungtrue\fi
\newcommand{\lsgfarbe}{\color{red!75!black}}

% ---------- Kopf- und Fußzeile (einheitliche Form) ----------
\newcommand{\lehrkraft}{Hau}
\newcommand{\themenbereich}{1. Formale Sprachen}
\newcommand{\abnummer}{}
\pagestyle{fancy}\fancyhf{}
\renewcommand{\headrulewidth}{0.4pt}
\fancyhead[L]{\small Klasse 13\\Datum:}
\fancyhead[C]{\small Themenbereich:\\\themenbereich}
\fancyhead[R]{\small \lehrkraft\ifloesung\\{\lsgfarbe\bfseries Lösung}\fi}
\fancyfoot[C]{\scriptsize Informatik 13 gA\,·\,Blatt \abnummer\ifloesung\ (Lösung)\fi\,·\,Seite \thepage}
\newcommand{\abtitel}[2]{\renewcommand{\abnummer}{#1}\begin{center}\Large\bfseries #1\quad\underline{#2}\end{center}\vspace{-2mm}}

% ---------- Boxen ----------
\newenvironment{aufgabe}[2][]{\begin{tcolorbox}[enhanced,breakable,colback=white,colframe=black,boxrule=0.7pt,arc=3mm,left=2mm,right=2mm,top=1mm,bottom=1.2mm,before skip=2mm,after skip=2mm]\textbf{Aufgabe #2)}\ifstrempty{#1}{}{ \textit{#1}}\par\vspace{0.3mm}}{\end{tcolorbox}}
\newtcolorbox{merke}[1][Merke]{enhanced,breakable,colback=green!5,colframe=green!40!black,boxrule=0.8pt,arc=2mm,left=2mm,right=2mm,top=1mm,bottom=1.2mm,before skip=2mm,after skip=2mm,title={#1},fonttitle=\bfseries,coltitle=white,toptitle=0.3mm,bottomtitle=0.3mm}
\newtcolorbox{hinweis}{enhanced,breakable,colback=yellow!12,colframe=orange!70!black,boxrule=0.6pt,arc=2mm,left=2mm,right=2mm,top=1mm,bottom=1mm,before skip=2mm,after skip=2mm}

% ---------- Lücken, Linien, Karos ----------
\newcommand{\luecke}[1][3.5cm]{\rule[-0.5ex]{#1}{0.4pt}}
\newcommand{\linien}[1]{\par\vspace{1.5mm}\foreach \i in {1,...,#1}{\noindent\rule[0pt]{\linewidth}{0.3pt}\par\vspace{5.5mm}}\vspace{-3mm}}
\newcommand{\code}[1]{\texttt{#1}}
\newcommand{\eps}{$\varepsilon$}

% Lücke mit Lösung: \lsg[Mindestbreite]{Lösungstext}
%  Die Lösung steht (ohne eigene Höhe/Tiefe) auf derselben Linie wie im Blatt: Das Layout verschiebt sich nicht,
%  solange der Text nicht breiter als die Mindestbreite ist.
\newlength{\lsgbreite}
\newcommand{\lsg}[2][3.5cm]{\leavevmode\ifloesung\settowidth{\lsgbreite}{\,#2\,}\ifdim\lsgbreite<#1\setlength{\lsgbreite}{#1}\fi\rlap{\smash{\raisebox{0.8pt}{\makebox[\lsgbreite][l]{\lsgfarbe\,#2}}}}\luecke[\lsgbreite]\else\luecke[#1]\fi}
% Schreiblinien mit Lösung: \lsglinien{Anzahl}{Lösungstext} (gleiche Höhe wie \linien:
%  je Linie \baselineskip + 5,5 mm + \parskip, ausgemessen)
\newlength{\lsglinienhoehe}
\newcommand{\lsglinien}[2]{\ifloesung\par\vspace{1.5mm}\setlength{\lsglinienhoehe}{\dimexpr #1\dimexpr\baselineskip+5.5mm+\parskip\relax-16.25pt\relax}\noindent\begin{minipage}[t][\lsglinienhoehe][t]{\linewidth}\lsgfarbe #2\end{minipage}\par\vspace{-1.5mm}\else\linien{#1}\fi}
% Beliebiger Inhalt: \lsgoder{nur in der Lösung}{nur im Blatt}
\newcommand{\lsgoder}[2]{\ifloesung\leavevmode#1\else#2\fi}
% Text nur in der Lösung sichtbar, belegt im Blatt aber denselben Platz
\newcommand{\lsgtext}[1]{\ifloesung{\lsgfarbe#1}\else\phantom{#1}\fi}
% Ankreuzfeld: \kreuz{1} = angekreuzt (nur in der Lösung), \kreuz{0} = leer
\newcommand{\kreuz}[1]{\leavevmode\ifloesung\ifnum#1=1 {\lsgfarbe$\boxtimes$}\else$\square$\fi\else$\square$\fi}

% ---------- Syntaxdiagramme (Maße in mm) ----------
%  Terminal rund (orange), Nichtterminal eckig (grün), Linien mit Pfeilen:
%  \draw[sdl,->] (0,0) -- ++(6,0) node[nt,anchor=west] (a) {Augen};   (Pfeil endet am Kasten)
\tikzset{
  sd/.style={x=1mm,y=1mm,>={Stealth[length=1.6mm,width=1.3mm]},font=\small},
  sdl/.style={line width=0.55pt,rounded corners=1.2mm},
  te/.style={draw,rounded corners=2.1mm,minimum height=4.2mm,minimum width=4.2mm,inner sep=1pt,font=\ttfamily\small,fill=orange!18},
  nt/.style={draw,minimum height=4.2mm,inner sep=1.6pt,font=\small,fill=green!12},
  sdname/.style={anchor=south west,font=\small\bfseries,inner sep=0pt},
  punkte/.style={font=\small,inner sep=0pt},
  % Ableitungsbaum
  baum/.style={font=\small,edge from parent/.style={draw,line width=0.5pt,-{Stealth[length=1.4mm]}},level distance=8.5mm},
}
% Lösungszeichnung: im Blatt unsichtbar, belegt aber denselben Platz
\ifloesung
  \tikzset{lsgbild/.style={red!75!black,te/.append style={fill=red!5},nt/.append style={fill=red!5}}}
\else
  \tikzset{lsgbild/.style={draw opacity=0,fill opacity=0,text opacity=0}}
\fi
% Karofläche, in der Lösung mit roter Zeichnung: \karolsg{Anzahl Karozeilen}{TikZ-Code}
% Koordinaten der Zeichnung in mm, Ursprung links oben, y nach unten negativ.
\newcommand{\karolsg}[2]{\par\vspace{1mm}\noindent\begin{tikzpicture}
  \draw[step=5mm,gray!45,thin] (0,0) grid (\linewidth-1pt,#1*5mm);
  \useasboundingbox (0,0) rectangle (\linewidth-1pt,#1*5mm);
  \ifloesung\begin{scope}[sd,shift={(0,#1*5mm)},lsgbild]#2\end{scope}\fi
\end{tikzpicture}\par}

% ---------- Bausteine für Syntaxdiagramme (innerhalb tikzpicture[sd]) ----------
%  Jeder Baustein beginnt an einem Punkt (Knoten-Anker oder Koordinate) und
%  legt die Koordinate <name>-o an, an der es weitergeht.
% Aneinanderreihung: \sdnext{name}{von}{te|nt}{Text}   (weiter bei name.east)
\newcommand{\sdnext}[4]{\draw[sdl,->] (#2) -- ++(4,0) node[#3,anchor=west] (#1) {#4};\coordinate (#1-o) at (#1.east);}
% Option: \sdopt{name}{von}{te|nt}{Text}
\newcommand{\sdopt}[4]{\path (#2) ++(4,-5.5) node[#3,anchor=west] (#1) {#4};
  \coordinate (#1-x) at ($(#1.east)+(3.5,0)$); \coordinate (#1-o) at (#1-x |- #2);
  \draw[sdl,->] (#2) |- (#1.west); \draw[sdl,->] (#1.east) -| (#1-o); \draw[sdl] (#2) -- (#1-o);}
% Wiederholung, mindestens einmal: \sdwdh{name}{von}{te|nt}{Text}
\newcommand{\sdwdh}[4]{\draw[sdl,->] (#2) -- ++(6,0) node[#3,anchor=west] (#1) {#4};
  \coordinate (#1-o) at ($(#1.east)+(4,0)$);
  \draw[sdl] (#1.east) -- (#1-o); \draw[sdl,->] (#1-o) -- ++(0,-5.5) -| ($(#2)+(2.5,0)$);}
% Wiederholung, beliebig oft (auch keinmal): \sdwdhn{name}{von}{te|nt}{Text}
\newcommand{\sdwdhn}[4]{\path (#2) ++(6,-5.5) node[#3,anchor=west] (#1) {#4};
  \coordinate (#1-x) at ($(#1.east)+(4,0)$); \coordinate (#1-o) at (#1-x |- #2);
  \draw[sdl] (#2) -- (#1-o); \draw[sdl,->] (#1-o) |- (#1.east); \draw[sdl,->] (#1.west) -| ($(#2)+(2.5,0)$);}
% Alternative: \sdalt{name}{von}{Breite bis zum Zusammenführen}{Zeilenabstand}{te/a, nt/Ziffer, pk/, ...}
%  pk/ setzt senkrechte Pünktchen (…) statt eines Zweigs
\newcommand{\sdalt}[5]{\coordinate (#1-o) at ($(#2)+(#3,0)$);
  \foreach \typ/\txt [count=\i from 0] in {#5}{%
    \ifdefstring{\typ}{pk}{\node[punkte] at ($(#2)+(6.1,-\i*#4+0.4)$) {$\vdots$};}{%
      \ifdefstring{\typ}{te}{\tikzset{sdakt/.style={te}}}{\tikzset{sdakt/.style={nt}}}%
      \node[sdakt,anchor=west] (#1\i) at ($(#2)+(4,-\i*#4)$) {\txt};
      \ifnum\i=0\relax \draw[sdl,->] (#2) -- (#1\i.west); \draw[sdl] (#1\i.east) -- (#1-o);
      \else \draw[sdl,->] (#2) |- (#1\i.west); \draw[sdl,->] (#1\i.east) -| (#1-o);\fi}}}
% Ende eines Diagramms
\newcommand{\sdende}[1]{\draw[sdl,->] (#1) -- ++(5,0);}

% ---------- Die Smiley-Grammatik (Blatt 1.3 und 1.4) ----------
\newcommand{\sdSmileyfolge}{\begin{tikzpicture}[sd,baseline=0mm]
  \node[sdname] at (0,3.3) {Smileyfolge}; \coordinate (s) at (0,0);
  \sdwdh{sm}{s}{nt}{Smiley}\sdende{sm-o}\end{tikzpicture}}
\newcommand{\sdSmiley}{\begin{tikzpicture}[sd,baseline=0mm]
  \node[sdname] at (0,3.3) {Smiley}; \coordinate (s) at (0,0);
  \sdnext{au}{s}{nt}{Augen}\sdnext{mi}{au-o}{nt}{Mitte}\sdnext{mu}{mi-o}{nt}{Mund}\sdende{mu-o}\end{tikzpicture}}
\newcommand{\sdAugen}{\begin{tikzpicture}[sd,baseline=0mm]
  \node[sdname] at (0,3.3) {Augen}; \coordinate (s) at (0,0);
  \sdalt{a}{s}{12}{5.5}{te/:, te/;}\sdende{a-o}\end{tikzpicture}}
\newcommand{\sdMitte}{\begin{tikzpicture}[sd,baseline=0mm]
  \node[sdname] at (0,3.3) {Mitte}; \coordinate (s) at (0,0);
  \sdopt{n}{s}{te}{-}\sdende{n-o}\end{tikzpicture}}
\newcommand{\sdMund}{\begin{tikzpicture}[sd,baseline=0mm]
  \node[sdname] at (0,3.3) {Mund}; \coordinate (s) at (0,0);
  \sdalt{m}{s}{12}{5.5}{te/), te/{(}, te/D, te/P, te/*}\sdende{m-o}\end{tikzpicture}}

%  Übersetzen mit XeLaTeX, LuaLaTeX, pdfLaTeX oder Tectonic.
%
%  Blatt:   tectonic ab-1-2-formale-sprachen.tex
%  Lösung:  tectonic ab-1-2-formale-sprachen-lsg.tex
%           (Datei enthält nur: \def\mitloesung{}\documentclass[11pt,a4paper]{article}
% =====================================================================
%  Gemeinsame Vorlage der Arbeitsblätter – 1. Formale Sprachen (Q13)
%  (gleiches Layout wie informatik/13/material/grundwissen/ab-vorlage.tex
%  und informatik/12/rekursive-datenstrukturen/arbeitsblaetter/ab-vorlage.tex)
%  Einbinden in jedem Blatt direkt nach \documentclass: \input{ab-vorlage}
%  Übersetzen mit XeLaTeX, LuaLaTeX, pdfLaTeX oder Tectonic.
%
%  Blatt:   tectonic ab-1-2-formale-sprachen.tex
%  Lösung:  tectonic ab-1-2-formale-sprachen-lsg.tex
%           (Datei enthält nur: \def\mitloesung{}\input{ab-1-2-formale-sprachen};
%            füllt alle Lücken rot aus; siehe \lsg, \lsglinien, \lsgoder, lsgbild, \karolsg)
% =====================================================================
\usepackage[a4paper,left=16mm,right=16mm,top=23mm,bottom=13mm,headheight=12mm,headsep=3mm,footskip=6mm]{geometry}
\usepackage{iftex}
\ifPDFTeX
  \usepackage[T1]{fontenc}
  \usepackage[utf8]{inputenc}
  \usepackage{helvet}
  \usepackage[scaled=0.86]{beramono}
  \renewcommand{\familydefault}{\sfdefault}
\else
  \usepackage{fontspec}
  \setmainfont{texgyreheros}[Extension=.otf,UprightFont=*-regular,BoldFont=*-bold,ItalicFont=*-italic,BoldItalicFont=*-bolditalic]
  \setmonofont{DejaVuSansMono}[Extension=.ttf,UprightFont=*,BoldFont=*-Bold,ItalicFont=*-Oblique,BoldItalicFont=*-BoldOblique,Scale=0.86]
\fi
\usepackage[ngerman,shorthands=off]{babel}
\usepackage{amssymb}
\usepackage{xcolor,graphicx,tabularx,array,enumitem,fancyhdr,tikz,multicol,calc,etoolbox,pifont}
\usepackage[most]{tcolorbox}
\usetikzlibrary{arrows.meta,positioning,calc}
\usepackage[hidelinks]{hyperref}
\setlength{\parindent}{0pt}
\setlength{\parskip}{3pt}
\setlist{nosep,leftmargin=*}

% ---------- Lösungsschalter ----------
\newif\ifloesung
\ifdefined\mitloesung\loesungtrue\fi
\newcommand{\lsgfarbe}{\color{red!75!black}}

% ---------- Kopf- und Fußzeile (einheitliche Form) ----------
\newcommand{\lehrkraft}{Hau}
\newcommand{\themenbereich}{1. Formale Sprachen}
\newcommand{\abnummer}{}
\pagestyle{fancy}\fancyhf{}
\renewcommand{\headrulewidth}{0.4pt}
\fancyhead[L]{\small Klasse 13\\Datum:}
\fancyhead[C]{\small Themenbereich:\\\themenbereich}
\fancyhead[R]{\small \lehrkraft\ifloesung\\{\lsgfarbe\bfseries Lösung}\fi}
\fancyfoot[C]{\scriptsize Informatik 13 gA\,·\,Blatt \abnummer\ifloesung\ (Lösung)\fi\,·\,Seite \thepage}
\newcommand{\abtitel}[2]{\renewcommand{\abnummer}{#1}\begin{center}\Large\bfseries #1\quad\underline{#2}\end{center}\vspace{-2mm}}

% ---------- Boxen ----------
\newenvironment{aufgabe}[2][]{\begin{tcolorbox}[enhanced,breakable,colback=white,colframe=black,boxrule=0.7pt,arc=3mm,left=2mm,right=2mm,top=1mm,bottom=1.2mm,before skip=2mm,after skip=2mm]\textbf{Aufgabe #2)}\ifstrempty{#1}{}{ \textit{#1}}\par\vspace{0.3mm}}{\end{tcolorbox}}
\newtcolorbox{merke}[1][Merke]{enhanced,breakable,colback=green!5,colframe=green!40!black,boxrule=0.8pt,arc=2mm,left=2mm,right=2mm,top=1mm,bottom=1.2mm,before skip=2mm,after skip=2mm,title={#1},fonttitle=\bfseries,coltitle=white,toptitle=0.3mm,bottomtitle=0.3mm}
\newtcolorbox{hinweis}{enhanced,breakable,colback=yellow!12,colframe=orange!70!black,boxrule=0.6pt,arc=2mm,left=2mm,right=2mm,top=1mm,bottom=1mm,before skip=2mm,after skip=2mm}

% ---------- Lücken, Linien, Karos ----------
\newcommand{\luecke}[1][3.5cm]{\rule[-0.5ex]{#1}{0.4pt}}
\newcommand{\linien}[1]{\par\vspace{1.5mm}\foreach \i in {1,...,#1}{\noindent\rule[0pt]{\linewidth}{0.3pt}\par\vspace{5.5mm}}\vspace{-3mm}}
\newcommand{\code}[1]{\texttt{#1}}
\newcommand{\eps}{$\varepsilon$}

% Lücke mit Lösung: \lsg[Mindestbreite]{Lösungstext}
%  Die Lösung steht (ohne eigene Höhe/Tiefe) auf derselben Linie wie im Blatt: Das Layout verschiebt sich nicht,
%  solange der Text nicht breiter als die Mindestbreite ist.
\newlength{\lsgbreite}
\newcommand{\lsg}[2][3.5cm]{\leavevmode\ifloesung\settowidth{\lsgbreite}{\,#2\,}\ifdim\lsgbreite<#1\setlength{\lsgbreite}{#1}\fi\rlap{\smash{\raisebox{0.8pt}{\makebox[\lsgbreite][l]{\lsgfarbe\,#2}}}}\luecke[\lsgbreite]\else\luecke[#1]\fi}
% Schreiblinien mit Lösung: \lsglinien{Anzahl}{Lösungstext} (gleiche Höhe wie \linien:
%  je Linie \baselineskip + 5,5 mm + \parskip, ausgemessen)
\newlength{\lsglinienhoehe}
\newcommand{\lsglinien}[2]{\ifloesung\par\vspace{1.5mm}\setlength{\lsglinienhoehe}{\dimexpr #1\dimexpr\baselineskip+5.5mm+\parskip\relax-16.25pt\relax}\noindent\begin{minipage}[t][\lsglinienhoehe][t]{\linewidth}\lsgfarbe #2\end{minipage}\par\vspace{-1.5mm}\else\linien{#1}\fi}
% Beliebiger Inhalt: \lsgoder{nur in der Lösung}{nur im Blatt}
\newcommand{\lsgoder}[2]{\ifloesung\leavevmode#1\else#2\fi}
% Text nur in der Lösung sichtbar, belegt im Blatt aber denselben Platz
\newcommand{\lsgtext}[1]{\ifloesung{\lsgfarbe#1}\else\phantom{#1}\fi}
% Ankreuzfeld: \kreuz{1} = angekreuzt (nur in der Lösung), \kreuz{0} = leer
\newcommand{\kreuz}[1]{\leavevmode\ifloesung\ifnum#1=1 {\lsgfarbe$\boxtimes$}\else$\square$\fi\else$\square$\fi}

% ---------- Syntaxdiagramme (Maße in mm) ----------
%  Terminal rund (orange), Nichtterminal eckig (grün), Linien mit Pfeilen:
%  \draw[sdl,->] (0,0) -- ++(6,0) node[nt,anchor=west] (a) {Augen};   (Pfeil endet am Kasten)
\tikzset{
  sd/.style={x=1mm,y=1mm,>={Stealth[length=1.6mm,width=1.3mm]},font=\small},
  sdl/.style={line width=0.55pt,rounded corners=1.2mm},
  te/.style={draw,rounded corners=2.1mm,minimum height=4.2mm,minimum width=4.2mm,inner sep=1pt,font=\ttfamily\small,fill=orange!18},
  nt/.style={draw,minimum height=4.2mm,inner sep=1.6pt,font=\small,fill=green!12},
  sdname/.style={anchor=south west,font=\small\bfseries,inner sep=0pt},
  punkte/.style={font=\small,inner sep=0pt},
  % Ableitungsbaum
  baum/.style={font=\small,edge from parent/.style={draw,line width=0.5pt,-{Stealth[length=1.4mm]}},level distance=8.5mm},
}
% Lösungszeichnung: im Blatt unsichtbar, belegt aber denselben Platz
\ifloesung
  \tikzset{lsgbild/.style={red!75!black,te/.append style={fill=red!5},nt/.append style={fill=red!5}}}
\else
  \tikzset{lsgbild/.style={draw opacity=0,fill opacity=0,text opacity=0}}
\fi
% Karofläche, in der Lösung mit roter Zeichnung: \karolsg{Anzahl Karozeilen}{TikZ-Code}
% Koordinaten der Zeichnung in mm, Ursprung links oben, y nach unten negativ.
\newcommand{\karolsg}[2]{\par\vspace{1mm}\noindent\begin{tikzpicture}
  \draw[step=5mm,gray!45,thin] (0,0) grid (\linewidth-1pt,#1*5mm);
  \useasboundingbox (0,0) rectangle (\linewidth-1pt,#1*5mm);
  \ifloesung\begin{scope}[sd,shift={(0,#1*5mm)},lsgbild]#2\end{scope}\fi
\end{tikzpicture}\par}

% ---------- Bausteine für Syntaxdiagramme (innerhalb tikzpicture[sd]) ----------
%  Jeder Baustein beginnt an einem Punkt (Knoten-Anker oder Koordinate) und
%  legt die Koordinate <name>-o an, an der es weitergeht.
% Aneinanderreihung: \sdnext{name}{von}{te|nt}{Text}   (weiter bei name.east)
\newcommand{\sdnext}[4]{\draw[sdl,->] (#2) -- ++(4,0) node[#3,anchor=west] (#1) {#4};\coordinate (#1-o) at (#1.east);}
% Option: \sdopt{name}{von}{te|nt}{Text}
\newcommand{\sdopt}[4]{\path (#2) ++(4,-5.5) node[#3,anchor=west] (#1) {#4};
  \coordinate (#1-x) at ($(#1.east)+(3.5,0)$); \coordinate (#1-o) at (#1-x |- #2);
  \draw[sdl,->] (#2) |- (#1.west); \draw[sdl,->] (#1.east) -| (#1-o); \draw[sdl] (#2) -- (#1-o);}
% Wiederholung, mindestens einmal: \sdwdh{name}{von}{te|nt}{Text}
\newcommand{\sdwdh}[4]{\draw[sdl,->] (#2) -- ++(6,0) node[#3,anchor=west] (#1) {#4};
  \coordinate (#1-o) at ($(#1.east)+(4,0)$);
  \draw[sdl] (#1.east) -- (#1-o); \draw[sdl,->] (#1-o) -- ++(0,-5.5) -| ($(#2)+(2.5,0)$);}
% Wiederholung, beliebig oft (auch keinmal): \sdwdhn{name}{von}{te|nt}{Text}
\newcommand{\sdwdhn}[4]{\path (#2) ++(6,-5.5) node[#3,anchor=west] (#1) {#4};
  \coordinate (#1-x) at ($(#1.east)+(4,0)$); \coordinate (#1-o) at (#1-x |- #2);
  \draw[sdl] (#2) -- (#1-o); \draw[sdl,->] (#1-o) |- (#1.east); \draw[sdl,->] (#1.west) -| ($(#2)+(2.5,0)$);}
% Alternative: \sdalt{name}{von}{Breite bis zum Zusammenführen}{Zeilenabstand}{te/a, nt/Ziffer, pk/, ...}
%  pk/ setzt senkrechte Pünktchen (…) statt eines Zweigs
\newcommand{\sdalt}[5]{\coordinate (#1-o) at ($(#2)+(#3,0)$);
  \foreach \typ/\txt [count=\i from 0] in {#5}{%
    \ifdefstring{\typ}{pk}{\node[punkte] at ($(#2)+(6.1,-\i*#4+0.4)$) {$\vdots$};}{%
      \ifdefstring{\typ}{te}{\tikzset{sdakt/.style={te}}}{\tikzset{sdakt/.style={nt}}}%
      \node[sdakt,anchor=west] (#1\i) at ($(#2)+(4,-\i*#4)$) {\txt};
      \ifnum\i=0\relax \draw[sdl,->] (#2) -- (#1\i.west); \draw[sdl] (#1\i.east) -- (#1-o);
      \else \draw[sdl,->] (#2) |- (#1\i.west); \draw[sdl,->] (#1\i.east) -| (#1-o);\fi}}}
% Ende eines Diagramms
\newcommand{\sdende}[1]{\draw[sdl,->] (#1) -- ++(5,0);}

% ---------- Die Smiley-Grammatik (Blatt 1.3 und 1.4) ----------
\newcommand{\sdSmileyfolge}{\begin{tikzpicture}[sd,baseline=0mm]
  \node[sdname] at (0,3.3) {Smileyfolge}; \coordinate (s) at (0,0);
  \sdwdh{sm}{s}{nt}{Smiley}\sdende{sm-o}\end{tikzpicture}}
\newcommand{\sdSmiley}{\begin{tikzpicture}[sd,baseline=0mm]
  \node[sdname] at (0,3.3) {Smiley}; \coordinate (s) at (0,0);
  \sdnext{au}{s}{nt}{Augen}\sdnext{mi}{au-o}{nt}{Mitte}\sdnext{mu}{mi-o}{nt}{Mund}\sdende{mu-o}\end{tikzpicture}}
\newcommand{\sdAugen}{\begin{tikzpicture}[sd,baseline=0mm]
  \node[sdname] at (0,3.3) {Augen}; \coordinate (s) at (0,0);
  \sdalt{a}{s}{12}{5.5}{te/:, te/;}\sdende{a-o}\end{tikzpicture}}
\newcommand{\sdMitte}{\begin{tikzpicture}[sd,baseline=0mm]
  \node[sdname] at (0,3.3) {Mitte}; \coordinate (s) at (0,0);
  \sdopt{n}{s}{te}{-}\sdende{n-o}\end{tikzpicture}}
\newcommand{\sdMund}{\begin{tikzpicture}[sd,baseline=0mm]
  \node[sdname] at (0,3.3) {Mund}; \coordinate (s) at (0,0);
  \sdalt{m}{s}{12}{5.5}{te/), te/{(}, te/D, te/P, te/*}\sdende{m-o}\end{tikzpicture}}

% Blatt:   tectonic ab-1-2-formale-sprachen.tex
% Lösung:  tectonic ab-1-2-formale-sprachen-lsg.tex
% Progression: Alphabet -> Wort/Länge -> Wortmenge A* (zählen: unendlich) -> Sprache L ⊆ A*
%              -> Transfer (Alltagssprachen, Syntax vs. Semantik) -> Smileys in Worten (mehrdeutig)
\newcommand{\zeile}{\rule[-2.2mm]{0pt}{7mm}}
\newcommand{\hoch}{\rule[-1.6mm]{0pt}{6mm}}
\begin{document}
\abtitel{1.2}{Formale Sprachen}

Ob \code{2 + * 5} oder \code{34.07.2026} korrekt gebildet sind, kann ein Computer nur prüfen, wenn \textbf{eindeutig} festgelegt ist, welche \emph{Zeichen} erlaubt sind und welche \emph{Zeichenfolgen} dazugehören.

\begin{merke}[Alphabet, Wort und Länge]
\begin{itemize}
\item Ein \textbf{Alphabet} $A$ ist eine \lsg[2.2cm]{endliche}, nichtleere Menge von \lsg[2cm]{Zeichen}.\\[0.5mm]
$A_1 = \{\code{0}, \code{1}\}$\qquad $A_2 = \{\code{a}, \code{b}, \code{c}, \dots, \code{z}\}$\qquad $A_3 = \{\code{:}, \code{;}, \code{-}, \code{)}, \code{(}, \code{D}, \code{P}, \code{*}\}$
\item Ein \textbf{Wort} (eine \textbf{Zeichenkette}) über $A$ ist eine endliche Folge von Zeichen aus $A$.\\
Das \textbf{leere Wort} \eps{} enthält \lsg[2.6cm]{kein Zeichen}.
\item Die \textbf{Länge} $|w|$ eines Wortes $w$ ist die \lsg[4cm]{Anzahl seiner Zeichen}. Beispiel: $|\code{0110}| = 4$, \ $|\varepsilon| = $ \lsg[0.8cm]{0}.
\end{itemize}
\end{merke}

\begin{aufgabe}{1}
\textbf{a)} Gib je drei Wörter unterschiedlicher Länge an und notiere ihre Länge.\\[1mm]
\renewcommand{\arraystretch}{1.1}
\begin{tabularx}{\linewidth}{|p{1.3cm}|X|X|X|}\hline
über $A_1$\zeile & $k_1 = \varepsilon$, \quad $|k_1| = 0$ & \lsgoder{\lsgfarbe $k_2 = \code{1}$, \quad $|k_2| = 1$}{} & \lsgoder{\lsgfarbe $k_3 = \code{0110}$, \quad $|k_3| = 4$}{}\\\hline
über $A_2$\zeile & \lsgoder{\lsgfarbe \code{x}, \quad Länge 1}{} & \lsgoder{\lsgfarbe \code{baum}, \quad Länge 4}{} & \lsgoder{\lsgfarbe \code{informatik}, \quad Länge 10}{}\\\hline
über $A_3$\zeile & \lsgoder{\lsgfarbe \code{;D}, \quad Länge 2}{} & \lsgoder{\lsgfarbe \code{:-)}, \quad Länge 3}{} & \lsgoder{\lsgfarbe \code{DD:*(}, \quad Länge 5}{}\\\hline
\end{tabularx}\\[1.2mm]
\textbf{b)} Ist \code{ab} ein Wort über $A_1$? Ist \code{:-)} ein Wort über $A_2$? \lsg[7.4cm]{Nein, beides nicht: Zeichen fehlen im Alphabet.}\\[1mm]
\textbf{c)} Es gibt genau ein Wort, das über \emph{jedem} Alphabet gebildet werden kann. Welches? Warum? \\[0.6mm]
\lsg[\linewidth]{Das leere Wort \eps: Es enthält kein Zeichen, also auch keines, das im Alphabet fehlen könnte.}\\[1mm]
\textbf{d)} Kreuze an, was richtig ist.\qquad
\kreuz{0} $\varepsilon \in A_1$\qquad \kreuz{1} $\code{0} \in A_1$\qquad \kreuz{1} \code{000} ist Wort über $A_1$\qquad \kreuz{0} $|\code{:-)}| = 1$
\end{aufgabe}

\begin{aufgabe}[Wie viele Wörter gibt es?]{2}
\textbf{a)} Schreibe \emph{alle} Wörter über $A_1$ der Länge 0 bis 3 auf und zähle sie. Zähle dann auch für $A_2$ und $A_3$.\\[1mm]
\renewcommand{\arraystretch}{1.1}
\begin{tabularx}{\linewidth}{|p{2.9cm}|p{0.85cm}|p{1.5cm}|p{2.8cm}|X|p{1.4cm}|}\hline
\textbf{Länge $n$} & \textbf{0} & \textbf{1} & \textbf{2} & \textbf{3} & \textbf{$n$}\\\hline
Wörter über $A_1$\rule[-6mm]{0pt}{11mm} & \lsgoder{\lsgfarbe \eps}{} & \lsgoder{\lsgfarbe \code{0}, \code{1}}{} & \lsgoder{\lsgfarbe \code{00}, \code{01}, \code{10}, \code{11}}{} & \lsgoder{\lsgfarbe \code{000}, \code{001}, \code{010}, \code{011}, \code{100}, \code{101}, \code{110}, \code{111}}{} & \\\hline
Anzahl über $A_1$\hoch & \lsgoder{\lsgfarbe 1}{} & \lsgoder{\lsgfarbe 2}{} & \lsgoder{\lsgfarbe 4}{} & \lsgoder{\lsgfarbe 8}{} & \lsgoder{\lsgfarbe $2^n$}{}\\\hline
Anzahl über $A_2$\hoch & \lsgoder{\lsgfarbe 1}{} & \lsgoder{\lsgfarbe 26}{} & \lsgoder{\lsgfarbe $26^2 = 676$}{} & \lsgoder{\lsgfarbe $26^3 = 17\,576$}{} & \lsgoder{\lsgfarbe $26^n$}{}\\\hline
Anzahl über $A_3$\hoch & \lsgoder{\lsgfarbe 1}{} & \lsgoder{\lsgfarbe 8}{} & \lsgoder{\lsgfarbe 64}{} & \lsgoder{\lsgfarbe 512}{} & \lsgoder{\lsgfarbe $8^n$}{}\\\hline
\end{tabularx}\\[1.2mm]
\textbf{b)} Wie viele Wörter gibt es \emph{insgesamt} über $A_1$? Begründe, obwohl $A_1$ nur zwei Zeichen hat.\\[0.6mm]
\lsg[\linewidth]{Unendlich viele: Zu jeder Länge $n$ gibt es $2^n$ Wörter, und die Länge ist nicht begrenzt.}
\end{aufgabe}

\begin{merke}[Wortmenge und formale Sprache]
\begin{itemize}
\item Die \textbf{Wortmenge} $A^*$ ist die Menge \lsg[2.6cm]{aller Wörter} über $A$, einschließlich \eps. Obwohl $A$ endlich ist, ist $A^*$ immer \lsg[2.4cm]{unendlich}.\\[0.5mm]
Beispiel: $A_1^* = \{$\,\lsg[9cm]{\eps, \code{0}, \code{1}, \code{00}, \code{01}, \code{10}, \code{11}, \code{000}, \code{001}, \dots}\,$\}$
\item Eine \textbf{formale Sprache} $L$ über $A$ ist eine \textbf{eindeutig festgelegte} Menge von Wörtern über $A$, also eine \lsg[2.4cm]{Teilmenge} von $A^*$: $L \subseteq A^*$. Eindeutig heißt: Für \emph{jedes} Wort aus $A^*$ steht fest, ob es zu $L$ gehört oder nicht.
\end{itemize}
\end{merke}

\newpage
\begin{aufgabe}[Sprachen über $A_2$]{3}
\begin{minipage}[t]{0.585\linewidth}
\vspace{0pt}
$L_1$ = alle Wörter aus $A_2^*$ mit Länge kleiner gleich 4\\
$L_2$ = alle Wörter aus $A_2^*$, die mit \code{b} beginnen\\
$L_3$ = alle Wörter aus $A_2^*$, die das Teilwort \code{inf} enthalten\\[1mm]
\textbf{a)} Kreuze an, zu welcher Sprache das Wort gehört.\\[1mm]
\textbf{b)} Welche der drei Sprachen sind endlich? \lsg[1.8cm]{nur $L_1$}\\[1mm]
\textbf{c)} Finde ein Wort, das zu $L_1$, $L_2$ \emph{und} $L_3$ gehört:\\[0.6mm] \lsg[3.2cm]{\code{binf}}\ \ Gibt es mehrere? \lsg[2.6cm]{nein, nur \code{binf}}
\end{minipage}\hfill
\begin{minipage}[t]{0.39\linewidth}
\vspace{0pt}
\renewcommand{\arraystretch}{1.12}
\begin{tabularx}{\linewidth}{|X|c|c|c|}\hline
\textbf{Wort} & $L_1$ & $L_2$ & $L_3$\\\hline
\code{baum} & \kreuz{1} & \kreuz{1} & \kreuz{0}\\\hline
\code{informatik} & \kreuz{0} & \kreuz{0} & \kreuz{1}\\\hline
\eps & \kreuz{1} & \kreuz{0} & \kreuz{0}\\\hline
\code{inf} & \kreuz{1} & \kreuz{0} & \kreuz{1}\\\hline
\code{binfo} & \kreuz{0} & \kreuz{1} & \kreuz{1}\\\hline
\end{tabularx}
\end{minipage}
\end{aufgabe}

\begin{aufgabe}[Formale Sprachen im Alltag]{4}
\textbf{a)} Ergänze die Tabelle.\\[1mm]
\renewcommand{\arraystretch}{1.1}
\begin{tabularx}{\linewidth}{|p{5.3cm}|X|p{2.4cm}|p{4.5cm}|}\hline
\textbf{Sprache} & \textbf{Alphabet} & \textbf{Wort aus $L$} & \textbf{Wort aus $A^*$, nicht aus $L$}\\\hline
Postleitzahlen\zeile & \lsgoder{\lsgfarbe \{0, 1, \dots, 9\}}{} & \code{83278} & \lsgoder{\lsgfarbe \code{824377} (6 Ziffern)}{}\\\hline
Binärzahlen ohne führende Null\zeile & \lsgoder{\lsgfarbe \{0, 1\}}{} & \lsgoder{\lsgfarbe \code{1011}}{} & \lsgoder{\lsgfarbe \code{0011}}{}\\\hline
Datum (TT.MM.JJJJ)\zeile & \lsgoder{\lsgfarbe \{0, 1, \dots, 9, .\}}{} & \lsgoder{\lsgfarbe \code{29.09.2026}}{} & \lsgoder{\lsgfarbe \code{29..092026}}{}\\\hline
\end{tabularx}\\[1.2mm]
\textbf{b)} Die Postleitzahl \code{99999} ist nicht vergeben. Gehört sie trotzdem zur Sprache „alle Wörter aus fünf Ziffern“? Was folgt daraus für Syntax und Semantik?
\lsglinien{2}{Ja, \code{99999} ist syntaktisch korrekt. Eine formale Sprache legt nur die \textbf{Syntax} fest; ob das Wort auch etwas bedeutet (\textbf{Semantik}), prüft sie nicht.}
\end{aufgabe}

\begin{aufgabe}[Die Sprache der Smileys]{5}
Im Chat werden Smileys aus den Zeichen von $A_3$ gebildet. Sie bilden eine Sprache $L_\text{Smiley} \subseteq A_3^*$.\\[1mm]
\textbf{a)} Kreuze an, welche Wörter du als Smiley erkennst.\\[1mm]
\kreuz{1} \code{:-)}\qquad \kreuz{1} \code{;D}\qquad \kreuz{1} \code{:-P}\qquad \kreuz{0} \code{)-:}\qquad \kreuz{0} \code{:-}\qquad \kreuz{1} \code{;*}\qquad \kreuz{0} \code{:--)}\qquad \kreuz{0} \code{DP}\qquad \kreuz{1} \code{:(}\\[1mm]
\textbf{b)} Beschreibe in Worten möglichst genau, welche Wörter aus $A_3^*$ zu $L_\text{Smiley}$ gehören.
\lsglinien{3}{Ein Smiley beginnt mit Augen (\code{:} oder \code{;}). Danach kann eine Nase \code{-} folgen, höchstens eine. Am Ende steht genau ein Mund: \code{)}, \code{(}, \code{D}, \code{P} oder \code{*}. Sonst folgt nichts.}
\textbf{c)} Vergleiche mit deinem Nachbarn: Gibt es Wörter, bei denen eure Beschreibungen verschieden entscheiden oder gar nicht entscheiden? Notiere ein Beispiel und woran es liegt.
\lsglinien{2}{z.\,B. \code{)-:} (gespiegelt?), \code{:--)} (wie viele Nasen?), \code{:-):-)} (mehrere Smileys?). Umgangssprache ist ungenau und mehrdeutig; man übersieht leicht Fälle.}
\end{aufgabe}

\begin{merke}[Erkenntnis]
Eine Beschreibung in Umgangssprache legt eine Sprache oft \lsg[2.8cm]{nicht eindeutig} fest. Nötig sind genaue \textbf{Regeln}, nach denen man für jedes Wort entscheiden kann, ob es zur Sprache gehört.
\end{merke}
\end{document}
;
%            füllt alle Lücken rot aus; siehe \lsg, \lsglinien, \lsgoder, lsgbild, \karolsg)
% =====================================================================
\usepackage[a4paper,left=16mm,right=16mm,top=23mm,bottom=13mm,headheight=12mm,headsep=3mm,footskip=6mm]{geometry}
\usepackage{iftex}
\ifPDFTeX
  \usepackage[T1]{fontenc}
  \usepackage[utf8]{inputenc}
  \usepackage{helvet}
  \usepackage[scaled=0.86]{beramono}
  \renewcommand{\familydefault}{\sfdefault}
\else
  \usepackage{fontspec}
  \setmainfont{texgyreheros}[Extension=.otf,UprightFont=*-regular,BoldFont=*-bold,ItalicFont=*-italic,BoldItalicFont=*-bolditalic]
  \setmonofont{DejaVuSansMono}[Extension=.ttf,UprightFont=*,BoldFont=*-Bold,ItalicFont=*-Oblique,BoldItalicFont=*-BoldOblique,Scale=0.86]
\fi
\usepackage[ngerman,shorthands=off]{babel}
\usepackage{amssymb}
\usepackage{xcolor,graphicx,tabularx,array,enumitem,fancyhdr,tikz,multicol,calc,etoolbox,pifont}
\usepackage[most]{tcolorbox}
\usetikzlibrary{arrows.meta,positioning,calc}
\usepackage[hidelinks]{hyperref}
\setlength{\parindent}{0pt}
\setlength{\parskip}{3pt}
\setlist{nosep,leftmargin=*}

% ---------- Lösungsschalter ----------
\newif\ifloesung
\ifdefined\mitloesung\loesungtrue\fi
\newcommand{\lsgfarbe}{\color{red!75!black}}

% ---------- Kopf- und Fußzeile (einheitliche Form) ----------
\newcommand{\lehrkraft}{Hau}
\newcommand{\themenbereich}{1. Formale Sprachen}
\newcommand{\abnummer}{}
\pagestyle{fancy}\fancyhf{}
\renewcommand{\headrulewidth}{0.4pt}
\fancyhead[L]{\small Klasse 13\\Datum:}
\fancyhead[C]{\small Themenbereich:\\\themenbereich}
\fancyhead[R]{\small \lehrkraft\ifloesung\\{\lsgfarbe\bfseries Lösung}\fi}
\fancyfoot[C]{\scriptsize Informatik 13 gA\,·\,Blatt \abnummer\ifloesung\ (Lösung)\fi\,·\,Seite \thepage}
\newcommand{\abtitel}[2]{\renewcommand{\abnummer}{#1}\begin{center}\Large\bfseries #1\quad\underline{#2}\end{center}\vspace{-2mm}}

% ---------- Boxen ----------
\newenvironment{aufgabe}[2][]{\begin{tcolorbox}[enhanced,breakable,colback=white,colframe=black,boxrule=0.7pt,arc=3mm,left=2mm,right=2mm,top=1mm,bottom=1.2mm,before skip=2mm,after skip=2mm]\textbf{Aufgabe #2)}\ifstrempty{#1}{}{ \textit{#1}}\par\vspace{0.3mm}}{\end{tcolorbox}}
\newtcolorbox{merke}[1][Merke]{enhanced,breakable,colback=green!5,colframe=green!40!black,boxrule=0.8pt,arc=2mm,left=2mm,right=2mm,top=1mm,bottom=1.2mm,before skip=2mm,after skip=2mm,title={#1},fonttitle=\bfseries,coltitle=white,toptitle=0.3mm,bottomtitle=0.3mm}
\newtcolorbox{hinweis}{enhanced,breakable,colback=yellow!12,colframe=orange!70!black,boxrule=0.6pt,arc=2mm,left=2mm,right=2mm,top=1mm,bottom=1mm,before skip=2mm,after skip=2mm}

% ---------- Lücken, Linien, Karos ----------
\newcommand{\luecke}[1][3.5cm]{\rule[-0.5ex]{#1}{0.4pt}}
\newcommand{\linien}[1]{\par\vspace{1.5mm}\foreach \i in {1,...,#1}{\noindent\rule[0pt]{\linewidth}{0.3pt}\par\vspace{5.5mm}}\vspace{-3mm}}
\newcommand{\code}[1]{\texttt{#1}}
\newcommand{\eps}{$\varepsilon$}

% Lücke mit Lösung: \lsg[Mindestbreite]{Lösungstext}
%  Die Lösung steht (ohne eigene Höhe/Tiefe) auf derselben Linie wie im Blatt: Das Layout verschiebt sich nicht,
%  solange der Text nicht breiter als die Mindestbreite ist.
\newlength{\lsgbreite}
\newcommand{\lsg}[2][3.5cm]{\leavevmode\ifloesung\settowidth{\lsgbreite}{\,#2\,}\ifdim\lsgbreite<#1\setlength{\lsgbreite}{#1}\fi\rlap{\smash{\raisebox{0.8pt}{\makebox[\lsgbreite][l]{\lsgfarbe\,#2}}}}\luecke[\lsgbreite]\else\luecke[#1]\fi}
% Schreiblinien mit Lösung: \lsglinien{Anzahl}{Lösungstext} (gleiche Höhe wie \linien:
%  je Linie \baselineskip + 5,5 mm + \parskip, ausgemessen)
\newlength{\lsglinienhoehe}
\newcommand{\lsglinien}[2]{\ifloesung\par\vspace{1.5mm}\setlength{\lsglinienhoehe}{\dimexpr #1\dimexpr\baselineskip+5.5mm+\parskip\relax-16.25pt\relax}\noindent\begin{minipage}[t][\lsglinienhoehe][t]{\linewidth}\lsgfarbe #2\end{minipage}\par\vspace{-1.5mm}\else\linien{#1}\fi}
% Beliebiger Inhalt: \lsgoder{nur in der Lösung}{nur im Blatt}
\newcommand{\lsgoder}[2]{\ifloesung\leavevmode#1\else#2\fi}
% Text nur in der Lösung sichtbar, belegt im Blatt aber denselben Platz
\newcommand{\lsgtext}[1]{\ifloesung{\lsgfarbe#1}\else\phantom{#1}\fi}
% Ankreuzfeld: \kreuz{1} = angekreuzt (nur in der Lösung), \kreuz{0} = leer
\newcommand{\kreuz}[1]{\leavevmode\ifloesung\ifnum#1=1 {\lsgfarbe$\boxtimes$}\else$\square$\fi\else$\square$\fi}

% ---------- Syntaxdiagramme (Maße in mm) ----------
%  Terminal rund (orange), Nichtterminal eckig (grün), Linien mit Pfeilen:
%  \draw[sdl,->] (0,0) -- ++(6,0) node[nt,anchor=west] (a) {Augen};   (Pfeil endet am Kasten)
\tikzset{
  sd/.style={x=1mm,y=1mm,>={Stealth[length=1.6mm,width=1.3mm]},font=\small},
  sdl/.style={line width=0.55pt,rounded corners=1.2mm},
  te/.style={draw,rounded corners=2.1mm,minimum height=4.2mm,minimum width=4.2mm,inner sep=1pt,font=\ttfamily\small,fill=orange!18},
  nt/.style={draw,minimum height=4.2mm,inner sep=1.6pt,font=\small,fill=green!12},
  sdname/.style={anchor=south west,font=\small\bfseries,inner sep=0pt},
  punkte/.style={font=\small,inner sep=0pt},
  % Ableitungsbaum
  baum/.style={font=\small,edge from parent/.style={draw,line width=0.5pt,-{Stealth[length=1.4mm]}},level distance=8.5mm},
}
% Lösungszeichnung: im Blatt unsichtbar, belegt aber denselben Platz
\ifloesung
  \tikzset{lsgbild/.style={red!75!black,te/.append style={fill=red!5},nt/.append style={fill=red!5}}}
\else
  \tikzset{lsgbild/.style={draw opacity=0,fill opacity=0,text opacity=0}}
\fi
% Karofläche, in der Lösung mit roter Zeichnung: \karolsg{Anzahl Karozeilen}{TikZ-Code}
% Koordinaten der Zeichnung in mm, Ursprung links oben, y nach unten negativ.
\newcommand{\karolsg}[2]{\par\vspace{1mm}\noindent\begin{tikzpicture}
  \draw[step=5mm,gray!45,thin] (0,0) grid (\linewidth-1pt,#1*5mm);
  \useasboundingbox (0,0) rectangle (\linewidth-1pt,#1*5mm);
  \ifloesung\begin{scope}[sd,shift={(0,#1*5mm)},lsgbild]#2\end{scope}\fi
\end{tikzpicture}\par}

% ---------- Bausteine für Syntaxdiagramme (innerhalb tikzpicture[sd]) ----------
%  Jeder Baustein beginnt an einem Punkt (Knoten-Anker oder Koordinate) und
%  legt die Koordinate <name>-o an, an der es weitergeht.
% Aneinanderreihung: \sdnext{name}{von}{te|nt}{Text}   (weiter bei name.east)
\newcommand{\sdnext}[4]{\draw[sdl,->] (#2) -- ++(4,0) node[#3,anchor=west] (#1) {#4};\coordinate (#1-o) at (#1.east);}
% Option: \sdopt{name}{von}{te|nt}{Text}
\newcommand{\sdopt}[4]{\path (#2) ++(4,-5.5) node[#3,anchor=west] (#1) {#4};
  \coordinate (#1-x) at ($(#1.east)+(3.5,0)$); \coordinate (#1-o) at (#1-x |- #2);
  \draw[sdl,->] (#2) |- (#1.west); \draw[sdl,->] (#1.east) -| (#1-o); \draw[sdl] (#2) -- (#1-o);}
% Wiederholung, mindestens einmal: \sdwdh{name}{von}{te|nt}{Text}
\newcommand{\sdwdh}[4]{\draw[sdl,->] (#2) -- ++(6,0) node[#3,anchor=west] (#1) {#4};
  \coordinate (#1-o) at ($(#1.east)+(4,0)$);
  \draw[sdl] (#1.east) -- (#1-o); \draw[sdl,->] (#1-o) -- ++(0,-5.5) -| ($(#2)+(2.5,0)$);}
% Wiederholung, beliebig oft (auch keinmal): \sdwdhn{name}{von}{te|nt}{Text}
\newcommand{\sdwdhn}[4]{\path (#2) ++(6,-5.5) node[#3,anchor=west] (#1) {#4};
  \coordinate (#1-x) at ($(#1.east)+(4,0)$); \coordinate (#1-o) at (#1-x |- #2);
  \draw[sdl] (#2) -- (#1-o); \draw[sdl,->] (#1-o) |- (#1.east); \draw[sdl,->] (#1.west) -| ($(#2)+(2.5,0)$);}
% Alternative: \sdalt{name}{von}{Breite bis zum Zusammenführen}{Zeilenabstand}{te/a, nt/Ziffer, pk/, ...}
%  pk/ setzt senkrechte Pünktchen (…) statt eines Zweigs
\newcommand{\sdalt}[5]{\coordinate (#1-o) at ($(#2)+(#3,0)$);
  \foreach \typ/\txt [count=\i from 0] in {#5}{%
    \ifdefstring{\typ}{pk}{\node[punkte] at ($(#2)+(6.1,-\i*#4+0.4)$) {$\vdots$};}{%
      \ifdefstring{\typ}{te}{\tikzset{sdakt/.style={te}}}{\tikzset{sdakt/.style={nt}}}%
      \node[sdakt,anchor=west] (#1\i) at ($(#2)+(4,-\i*#4)$) {\txt};
      \ifnum\i=0\relax \draw[sdl,->] (#2) -- (#1\i.west); \draw[sdl] (#1\i.east) -- (#1-o);
      \else \draw[sdl,->] (#2) |- (#1\i.west); \draw[sdl,->] (#1\i.east) -| (#1-o);\fi}}}
% Ende eines Diagramms
\newcommand{\sdende}[1]{\draw[sdl,->] (#1) -- ++(5,0);}

% ---------- Die Smiley-Grammatik (Blatt 1.3 und 1.4) ----------
\newcommand{\sdSmileyfolge}{\begin{tikzpicture}[sd,baseline=0mm]
  \node[sdname] at (0,3.3) {Smileyfolge}; \coordinate (s) at (0,0);
  \sdwdh{sm}{s}{nt}{Smiley}\sdende{sm-o}\end{tikzpicture}}
\newcommand{\sdSmiley}{\begin{tikzpicture}[sd,baseline=0mm]
  \node[sdname] at (0,3.3) {Smiley}; \coordinate (s) at (0,0);
  \sdnext{au}{s}{nt}{Augen}\sdnext{mi}{au-o}{nt}{Mitte}\sdnext{mu}{mi-o}{nt}{Mund}\sdende{mu-o}\end{tikzpicture}}
\newcommand{\sdAugen}{\begin{tikzpicture}[sd,baseline=0mm]
  \node[sdname] at (0,3.3) {Augen}; \coordinate (s) at (0,0);
  \sdalt{a}{s}{12}{5.5}{te/:, te/;}\sdende{a-o}\end{tikzpicture}}
\newcommand{\sdMitte}{\begin{tikzpicture}[sd,baseline=0mm]
  \node[sdname] at (0,3.3) {Mitte}; \coordinate (s) at (0,0);
  \sdopt{n}{s}{te}{-}\sdende{n-o}\end{tikzpicture}}
\newcommand{\sdMund}{\begin{tikzpicture}[sd,baseline=0mm]
  \node[sdname] at (0,3.3) {Mund}; \coordinate (s) at (0,0);
  \sdalt{m}{s}{12}{5.5}{te/), te/{(}, te/D, te/P, te/*}\sdende{m-o}\end{tikzpicture}}

% Blatt:   tectonic ab-1-2-formale-sprachen.tex
% Lösung:  tectonic ab-1-2-formale-sprachen-lsg.tex
% Progression: Alphabet -> Wort/Länge -> Wortmenge A* (zählen: unendlich) -> Sprache L ⊆ A*
%              -> Transfer (Alltagssprachen, Syntax vs. Semantik) -> Smileys in Worten (mehrdeutig)
\newcommand{\zeile}{\rule[-2.2mm]{0pt}{7mm}}
\newcommand{\hoch}{\rule[-1.6mm]{0pt}{6mm}}
\begin{document}
\abtitel{1.2}{Formale Sprachen}

Ob \code{2 + * 5} oder \code{34.07.2026} korrekt gebildet sind, kann ein Computer nur prüfen, wenn \textbf{eindeutig} festgelegt ist, welche \emph{Zeichen} erlaubt sind und welche \emph{Zeichenfolgen} dazugehören.

\begin{merke}[Alphabet, Wort und Länge]
\begin{itemize}
\item Ein \textbf{Alphabet} $A$ ist eine \lsg[2.2cm]{endliche}, nichtleere Menge von \lsg[2cm]{Zeichen}.\\[0.5mm]
$A_1 = \{\code{0}, \code{1}\}$\qquad $A_2 = \{\code{a}, \code{b}, \code{c}, \dots, \code{z}\}$\qquad $A_3 = \{\code{:}, \code{;}, \code{-}, \code{)}, \code{(}, \code{D}, \code{P}, \code{*}\}$
\item Ein \textbf{Wort} (eine \textbf{Zeichenkette}) über $A$ ist eine endliche Folge von Zeichen aus $A$.\\
Das \textbf{leere Wort} \eps{} enthält \lsg[2.6cm]{kein Zeichen}.
\item Die \textbf{Länge} $|w|$ eines Wortes $w$ ist die \lsg[4cm]{Anzahl seiner Zeichen}. Beispiel: $|\code{0110}| = 4$, \ $|\varepsilon| = $ \lsg[0.8cm]{0}.
\end{itemize}
\end{merke}

\begin{aufgabe}{1}
\textbf{a)} Gib je drei Wörter unterschiedlicher Länge an und notiere ihre Länge.\\[1mm]
\renewcommand{\arraystretch}{1.1}
\begin{tabularx}{\linewidth}{|p{1.3cm}|X|X|X|}\hline
über $A_1$\zeile & $k_1 = \varepsilon$, \quad $|k_1| = 0$ & \lsgoder{\lsgfarbe $k_2 = \code{1}$, \quad $|k_2| = 1$}{} & \lsgoder{\lsgfarbe $k_3 = \code{0110}$, \quad $|k_3| = 4$}{}\\\hline
über $A_2$\zeile & \lsgoder{\lsgfarbe \code{x}, \quad Länge 1}{} & \lsgoder{\lsgfarbe \code{baum}, \quad Länge 4}{} & \lsgoder{\lsgfarbe \code{informatik}, \quad Länge 10}{}\\\hline
über $A_3$\zeile & \lsgoder{\lsgfarbe \code{;D}, \quad Länge 2}{} & \lsgoder{\lsgfarbe \code{:-)}, \quad Länge 3}{} & \lsgoder{\lsgfarbe \code{DD:*(}, \quad Länge 5}{}\\\hline
\end{tabularx}\\[1.2mm]
\textbf{b)} Ist \code{ab} ein Wort über $A_1$? Ist \code{:-)} ein Wort über $A_2$? \lsg[7.4cm]{Nein, beides nicht: Zeichen fehlen im Alphabet.}\\[1mm]
\textbf{c)} Es gibt genau ein Wort, das über \emph{jedem} Alphabet gebildet werden kann. Welches? Warum? \\[0.6mm]
\lsg[\linewidth]{Das leere Wort \eps: Es enthält kein Zeichen, also auch keines, das im Alphabet fehlen könnte.}\\[1mm]
\textbf{d)} Kreuze an, was richtig ist.\qquad
\kreuz{0} $\varepsilon \in A_1$\qquad \kreuz{1} $\code{0} \in A_1$\qquad \kreuz{1} \code{000} ist Wort über $A_1$\qquad \kreuz{0} $|\code{:-)}| = 1$
\end{aufgabe}

\begin{aufgabe}[Wie viele Wörter gibt es?]{2}
\textbf{a)} Schreibe \emph{alle} Wörter über $A_1$ der Länge 0 bis 3 auf und zähle sie. Zähle dann auch für $A_2$ und $A_3$.\\[1mm]
\renewcommand{\arraystretch}{1.1}
\begin{tabularx}{\linewidth}{|p{2.9cm}|p{0.85cm}|p{1.5cm}|p{2.8cm}|X|p{1.4cm}|}\hline
\textbf{Länge $n$} & \textbf{0} & \textbf{1} & \textbf{2} & \textbf{3} & \textbf{$n$}\\\hline
Wörter über $A_1$\rule[-6mm]{0pt}{11mm} & \lsgoder{\lsgfarbe \eps}{} & \lsgoder{\lsgfarbe \code{0}, \code{1}}{} & \lsgoder{\lsgfarbe \code{00}, \code{01}, \code{10}, \code{11}}{} & \lsgoder{\lsgfarbe \code{000}, \code{001}, \code{010}, \code{011}, \code{100}, \code{101}, \code{110}, \code{111}}{} & \\\hline
Anzahl über $A_1$\hoch & \lsgoder{\lsgfarbe 1}{} & \lsgoder{\lsgfarbe 2}{} & \lsgoder{\lsgfarbe 4}{} & \lsgoder{\lsgfarbe 8}{} & \lsgoder{\lsgfarbe $2^n$}{}\\\hline
Anzahl über $A_2$\hoch & \lsgoder{\lsgfarbe 1}{} & \lsgoder{\lsgfarbe 26}{} & \lsgoder{\lsgfarbe $26^2 = 676$}{} & \lsgoder{\lsgfarbe $26^3 = 17\,576$}{} & \lsgoder{\lsgfarbe $26^n$}{}\\\hline
Anzahl über $A_3$\hoch & \lsgoder{\lsgfarbe 1}{} & \lsgoder{\lsgfarbe 8}{} & \lsgoder{\lsgfarbe 64}{} & \lsgoder{\lsgfarbe 512}{} & \lsgoder{\lsgfarbe $8^n$}{}\\\hline
\end{tabularx}\\[1.2mm]
\textbf{b)} Wie viele Wörter gibt es \emph{insgesamt} über $A_1$? Begründe, obwohl $A_1$ nur zwei Zeichen hat.\\[0.6mm]
\lsg[\linewidth]{Unendlich viele: Zu jeder Länge $n$ gibt es $2^n$ Wörter, und die Länge ist nicht begrenzt.}
\end{aufgabe}

\begin{merke}[Wortmenge und formale Sprache]
\begin{itemize}
\item Die \textbf{Wortmenge} $A^*$ ist die Menge \lsg[2.6cm]{aller Wörter} über $A$, einschließlich \eps. Obwohl $A$ endlich ist, ist $A^*$ immer \lsg[2.4cm]{unendlich}.\\[0.5mm]
Beispiel: $A_1^* = \{$\,\lsg[9cm]{\eps, \code{0}, \code{1}, \code{00}, \code{01}, \code{10}, \code{11}, \code{000}, \code{001}, \dots}\,$\}$
\item Eine \textbf{formale Sprache} $L$ über $A$ ist eine \textbf{eindeutig festgelegte} Menge von Wörtern über $A$, also eine \lsg[2.4cm]{Teilmenge} von $A^*$: $L \subseteq A^*$. Eindeutig heißt: Für \emph{jedes} Wort aus $A^*$ steht fest, ob es zu $L$ gehört oder nicht.
\end{itemize}
\end{merke}

\newpage
\begin{aufgabe}[Sprachen über $A_2$]{3}
\begin{minipage}[t]{0.585\linewidth}
\vspace{0pt}
$L_1$ = alle Wörter aus $A_2^*$ mit Länge kleiner gleich 4\\
$L_2$ = alle Wörter aus $A_2^*$, die mit \code{b} beginnen\\
$L_3$ = alle Wörter aus $A_2^*$, die das Teilwort \code{inf} enthalten\\[1mm]
\textbf{a)} Kreuze an, zu welcher Sprache das Wort gehört.\\[1mm]
\textbf{b)} Welche der drei Sprachen sind endlich? \lsg[1.8cm]{nur $L_1$}\\[1mm]
\textbf{c)} Finde ein Wort, das zu $L_1$, $L_2$ \emph{und} $L_3$ gehört:\\[0.6mm] \lsg[3.2cm]{\code{binf}}\ \ Gibt es mehrere? \lsg[2.6cm]{nein, nur \code{binf}}
\end{minipage}\hfill
\begin{minipage}[t]{0.39\linewidth}
\vspace{0pt}
\renewcommand{\arraystretch}{1.12}
\begin{tabularx}{\linewidth}{|X|c|c|c|}\hline
\textbf{Wort} & $L_1$ & $L_2$ & $L_3$\\\hline
\code{baum} & \kreuz{1} & \kreuz{1} & \kreuz{0}\\\hline
\code{informatik} & \kreuz{0} & \kreuz{0} & \kreuz{1}\\\hline
\eps & \kreuz{1} & \kreuz{0} & \kreuz{0}\\\hline
\code{inf} & \kreuz{1} & \kreuz{0} & \kreuz{1}\\\hline
\code{binfo} & \kreuz{0} & \kreuz{1} & \kreuz{1}\\\hline
\end{tabularx}
\end{minipage}
\end{aufgabe}

\begin{aufgabe}[Formale Sprachen im Alltag]{4}
\textbf{a)} Ergänze die Tabelle.\\[1mm]
\renewcommand{\arraystretch}{1.1}
\begin{tabularx}{\linewidth}{|p{5.3cm}|X|p{2.4cm}|p{4.5cm}|}\hline
\textbf{Sprache} & \textbf{Alphabet} & \textbf{Wort aus $L$} & \textbf{Wort aus $A^*$, nicht aus $L$}\\\hline
Postleitzahlen\zeile & \lsgoder{\lsgfarbe \{0, 1, \dots, 9\}}{} & \code{83278} & \lsgoder{\lsgfarbe \code{824377} (6 Ziffern)}{}\\\hline
Binärzahlen ohne führende Null\zeile & \lsgoder{\lsgfarbe \{0, 1\}}{} & \lsgoder{\lsgfarbe \code{1011}}{} & \lsgoder{\lsgfarbe \code{0011}}{}\\\hline
Datum (TT.MM.JJJJ)\zeile & \lsgoder{\lsgfarbe \{0, 1, \dots, 9, .\}}{} & \lsgoder{\lsgfarbe \code{29.09.2026}}{} & \lsgoder{\lsgfarbe \code{29..092026}}{}\\\hline
\end{tabularx}\\[1.2mm]
\textbf{b)} Die Postleitzahl \code{99999} ist nicht vergeben. Gehört sie trotzdem zur Sprache „alle Wörter aus fünf Ziffern“? Was folgt daraus für Syntax und Semantik?
\lsglinien{2}{Ja, \code{99999} ist syntaktisch korrekt. Eine formale Sprache legt nur die \textbf{Syntax} fest; ob das Wort auch etwas bedeutet (\textbf{Semantik}), prüft sie nicht.}
\end{aufgabe}

\begin{aufgabe}[Die Sprache der Smileys]{5}
Im Chat werden Smileys aus den Zeichen von $A_3$ gebildet. Sie bilden eine Sprache $L_\text{Smiley} \subseteq A_3^*$.\\[1mm]
\textbf{a)} Kreuze an, welche Wörter du als Smiley erkennst.\\[1mm]
\kreuz{1} \code{:-)}\qquad \kreuz{1} \code{;D}\qquad \kreuz{1} \code{:-P}\qquad \kreuz{0} \code{)-:}\qquad \kreuz{0} \code{:-}\qquad \kreuz{1} \code{;*}\qquad \kreuz{0} \code{:--)}\qquad \kreuz{0} \code{DP}\qquad \kreuz{1} \code{:(}\\[1mm]
\textbf{b)} Beschreibe in Worten möglichst genau, welche Wörter aus $A_3^*$ zu $L_\text{Smiley}$ gehören.
\lsglinien{3}{Ein Smiley beginnt mit Augen (\code{:} oder \code{;}). Danach kann eine Nase \code{-} folgen, höchstens eine. Am Ende steht genau ein Mund: \code{)}, \code{(}, \code{D}, \code{P} oder \code{*}. Sonst folgt nichts.}
\textbf{c)} Vergleiche mit deinem Nachbarn: Gibt es Wörter, bei denen eure Beschreibungen verschieden entscheiden oder gar nicht entscheiden? Notiere ein Beispiel und woran es liegt.
\lsglinien{2}{z.\,B. \code{)-:} (gespiegelt?), \code{:--)} (wie viele Nasen?), \code{:-):-)} (mehrere Smileys?). Umgangssprache ist ungenau und mehrdeutig; man übersieht leicht Fälle.}
\end{aufgabe}

\begin{merke}[Erkenntnis]
Eine Beschreibung in Umgangssprache legt eine Sprache oft \lsg[2.8cm]{nicht eindeutig} fest. Nötig sind genaue \textbf{Regeln}, nach denen man für jedes Wort entscheiden kann, ob es zur Sprache gehört.
\end{merke}
\end{document}
